\documentclass{article}
\usepackage[margin=1in]{geometry}
\usepackage{amsmath,amssymb,amsthm}
\usepackage{mathtools}
\usepackage{bm}
\usepackage{hyperref}
\usepackage{enumitem}
\usepackage{xcolor}
\usepackage{tcolorbox}
\tcbuselibrary{breakable}

\newtcolorbox{verdict}[1][]{enhanced,breakable,
  colback=blue!5,colframe=blue!60!black,
  fonttitle=\bfseries,title={#1},
  arc=2pt,boxrule=0.8pt,
  left=4pt,right=4pt,top=3pt,bottom=3pt,
  before skip=6pt,after skip=6pt}

\newtcolorbox{critical}[1][]{enhanced,breakable,
  colback=red!5,colframe=red!60!black,
  fonttitle=\bfseries,title={#1},
  arc=2pt,boxrule=0.8pt,
  left=4pt,right=4pt,top=3pt,bottom=3pt,
  before skip=6pt,after skip=6pt}

\newtcolorbox{suggestion}[1][]{enhanced,breakable,
  colback=green!5,colframe=green!60!black,
  fonttitle=\bfseries,title={#1},
  arc=2pt,boxrule=0.8pt,
  left=4pt,right=4pt,top=3pt,bottom=3pt,
  before skip=6pt,after skip=6pt}

\newtcolorbox{proofbox}[1][]{enhanced,breakable,
  colback=yellow!5,colframe=orange!70!black,
  fonttitle=\bfseries,title={#1},
  arc=2pt,boxrule=0.8pt,
  left=4pt,right=4pt,top=3pt,bottom=3pt,
  before skip=6pt,after skip=6pt}

\newtheorem{theorem}{Theorem}
\newtheorem{proposition}{Proposition}
\newtheorem{lemma}{Lemma}
\newtheorem{corollary}{Corollary}
\newtheorem{remark}{Remark}

\DeclareMathOperator{\tr}{tr}
\DeclareMathOperator{\col}{col}
\DeclareMathOperator{\diag}{diag}
\DeclareMathOperator{\op}{op}
\newcommand{\dR}{d^{\mathrm{R}}}
\newcommand{\dM}{d^{\mathrm{M}}}
\newcommand{\RR}{\mathbb{R}}

\title{Theoretical Evaluation of\\
\emph{Self-Consistent Riemannian Variational Autoencoders}\\[6pt]
\large Evaluation ID: Op46MX -- Version 1}

\author{Automated Theoretical Review}
\date{April 2026}

\begin{document}
\maketitle

\tableofcontents
\newpage

%% ====================================================================
\section{Meta-Review Summary}
%% ====================================================================

\begin{verdict}[Overall Assessment]
The SCR-VAE paper presents a novel and theoretically grounded approach to
learning isometric generative models of Riemannian data manifolds. The core
idea -- co-evolving the decoder's Riemannian metric and a data neighborhood
graph until self-consistency -- is creative and well-motivated. The fixed-point
characterization and convergence analysis under the Polyak-\L{}ojasiewicz
condition are technically solid. The curvature certificate via the ambient
closure vector is an elegant byproduct.

However, the evaluation identifies several \textbf{critical gaps}:
\begin{enumerate}[label=(\roman*)]
    \item A \textbf{fundamental data-generation error} in the torus experiment
      (standard embedded torus used instead of the intrinsically flat Clifford
      torus), which invalidates the primary negative-curvature test.
    \item An \textbf{internal contradiction} in the curvature proxy's
      theoretical prediction (Eq.~16 vs.\ Remark~8), leading to incorrect
      experimental comparisons.
    \item The \textbf{unresolved Assumption A4'} (metric approximation), which
      remains the weakest link in the isometry convergence theorem.
    \item A \textbf{missing edge KL term} in the theoretical loss (Eq.~3 of the
      paper) that is present in the implementation.
\end{enumerate}

\noindent
\textbf{Recommendation}: Major revision required before NeurIPS submission. The
core theory is sound but the experimental validation and certain theoretical
claims need correction.
\end{verdict}


%% ====================================================================
\section{Theory Review: NeurIPS Reviewer Assessment}
\label{sec:review}
%% ====================================================================

\subsection{Strengths}

\begin{enumerate}[label=\textbf{S\arabic*.}]
    \item \textbf{Novel and well-motivated framework.} The self-consistent
      iteration between decoder Riemannian metric and graph neighborhood
      structure is a genuine conceptual contribution. Unlike prior work on
      Riemannian VAEs (Arvanitidis et al.\ 2018), the paper co-evolves the
      graph with the metric rather than treating the graph as fixed.

    \item \textbf{Clean fixed-point theorem.} Theorem~1 is precisely stated and
      its three components (edge weight identity, distance approximation,
      principal tangent frame) follow naturally from the self-consistency
      definition. The proof of part~(b) via geodesic normal coordinates and
      the Bertrand-Puiseux expansion (Lemma~1) is rigorous and yields a
      tighter bound ($O(K_{\max} r^3)$) than originally claimed.

    \item \textbf{Honest convergence analysis.} The paper commendably admits
      that the E-step is NOT monotone (Remark~4), distinguishing the algorithm
      from standard EM. The convergence proof under the PL condition (Prop.~3)
      correctly handles the non-monotone perturbation and establishes geometric
      convergence to an $O(r_n^2)$ neighborhood.

    \item \textbf{Creative curvature certificate.} The ambient closure vector
      $c_{ijk}$ as a second-fundamental-form probe is elegant. Prop.~4
      correctly identifies that $c_{ijk}$ lies in the normal space to leading
      order, and the observation that the projected holonomy $W^+ c_{ijk}$
      destroys the curvature signal is an important negative result.

    \item \textbf{Transparent accounting.} Section~7 honestly lists what is
      proven, what is not, and what is open. This level of transparency is
      rare and valuable.
\end{enumerate}

\subsection{Weaknesses}

\begin{enumerate}[label=\textbf{W\arabic*.}]
    \item \textbf{Assumption A4' is unreduced (Critical).}
      The isometry convergence theorem (Theorem~2) depends on Assumption~A4'
      ($\|G^*(z_i) - g(z_i)\|_{\op} \leq C_G \delta_{\mathrm{rec}} + O(r_n^2)$),
      which the paper cannot derive from reconstruction accuracy alone.
      Remark~5 shows that the derivable bound has a $1/r_n$ factor:
      \[
          \|G^* - g\|_{\op} = O\!\left(\frac{\delta_{\mathrm{rec}}}{r_n} + r_n\right)
      \]
      This factor causes the Bernstein path accumulation in Step~4b to
      diverge as $n \to \infty$ for fixed $\delta_{\mathrm{rec}}$.
      The assumption requires either $\delta_{\mathrm{rec}} = O(r_n^2)$
      (very strong) or an independent Sobolev-type argument on the decoder.
      This is the single biggest theoretical gap.

    \item \textbf{Curvature proxy prediction is internally inconsistent (Critical).}
      Equation~(16) claims the scale-invariant prediction is
      $\hat{\kappa} \to 2/R = 2\sqrt{K}$. Remark~8 gives the correct formula
      for random embeddings: $\hat{\kappa} \to (2/R)\sqrt{3/G}$. For
      $R=1, G=50$: Eq.~(16) predicts $2.000$ while Remark~8 predicts
      $0.4899$. The experiments compare against $2.000$, which is incorrect
      for the random-embedding setting actually used. See Section~\ref{sec:kappa}
      for the detailed calculation.

    \item \textbf{Torus experiment uses the wrong manifold (Critical).}
      The data generation function \texttt{flat\_torus()} generates the
      standard embedded torus in $\RR^3$, which has non-zero Gaussian
      curvature $K(\phi) = \cos\phi / (r(R + r\cos\phi))$. This is NOT the
      flat torus. The geodesic evaluation uses the flat product metric
      $d = \sqrt{(R\,d\theta)^2 + (r\,d\phi)^2}$, creating a metric mismatch.
      See Section~\ref{sec:torus} for detailed analysis.

    \item \textbf{No global convergence guarantee.} The convergence is purely
      local (near a self-consistent fixed point). The basin of attraction is
      not characterized, and no condition is given for the Euclidean KNN
      initialization to land in the correct basin. For manifolds with multiple
      local minima, this is a significant limitation.

    \item \textbf{Mismatch between theoretical and implemented loss.}
      The paper's Eq.~(3) defines three loss terms:
      $\mathcal{L} = \mathcal{L}_{\text{recon}} + \beta\,\mathcal{L}_{\text{KL}}
      + \lambda\,\mathcal{L}_{\text{Riem}}$.
      The implementation adds a fourth term
      $\beta_e\,\mathcal{L}_{\text{edge\_KL}}$ not present in the theory.
      While this is a standard regularizer, its effect on the convergence
      analysis should be accounted for.
\end{enumerate}

\subsection{Questions for the Authors}

\begin{enumerate}[label=\textbf{Q\arabic*.}]
    \item Can Assumption~A4' be replaced by a condition on the decoder's
      Sobolev regularity (e.g., bounded $W^{1,\infty}$ norm of the weight
      matrices)?

    \item Which curvature proxy prediction is correct for the experimental
      setup: $2/R$ (Eq.~16) or $(2/R)\sqrt{3/G}$ (Remark~8)?

    \item Has the Clifford torus embedding in $\RR^4$ been tested? If so,
      does $\hat{\kappa}$ approach zero?

    \item What is the effect of the edge KL term on the convergence
      guarantees?

    \item For the sphere experiment, does the curvature proxy converge to the
      theoretical value as $N \to \infty$ and the graph becomes denser?
\end{enumerate}


%% ====================================================================
\section{Theory-Code Alignment Audit}
\label{sec:alignment}
%% ====================================================================

\subsection{Components Correctly Implemented}

\begin{enumerate}[label=\checkmark]
    \item \textbf{Stop-gradient on decoder for Riemannian targets.}
      In \texttt{log\_map.py}, the JVP computation uses
      \texttt{torch.no\_grad()} and returns \texttt{.detach()}.
      Additionally, the trainer passes \texttt{outputs["mu\_node"].detach()}
      as the \texttt{z\_mu} argument to the loss, ensuring no gradient flows
      from $\mathcal{L}_{\text{Riem}}$ to the encoder or decoder. This
      correctly implements Remark~1 of the paper.

    \item \textbf{Stiefel manifold constraint on $W$.}
      The \texttt{EdgeDecoder} initializes $W$ via \texttt{orthogonal\_init},
      and \texttt{retract\_to\_stiefel()} performs polar SVD retraction
      ($W \gets UV^\top$ where $W = U\Sigma V^\top$) after every gradient step.
      The Gram error is verified to be $0.000000$ at every iteration. This is
      theoretically correct and numerically superior to QR retraction.

    \item \textbf{Linear edge decoder without bias.}
      The \texttt{EdgeDecoder} computes $\hat{\ell}_{ij} = e_{ij} W^\top$
      with no bias term, matching the Eckart-Young requirement of Prop.~1.

    \item \textbf{Edge encoder operates on posterior means.}
      $\Delta z_{ij} = \mu_j - \mu_i$ (not reparameterized samples), matching
      the theory's deterministic specification.

    \item \textbf{Symmetrized Riemannian distances.}
      The \texttt{riemannian\_knn\_graph} computes both forward and backward
      log maps and averages: $\bar{w}_{ij} = (\|J_f(z_i)\Delta z_{ij}\| +
      \|J_f(z_j)\Delta z_{ji}\|)/2$. This is a sensible stabilization not
      discussed in the theory but consistent with the symmetric nature of
      Riemannian distance.

    \item \textbf{Distance clipping.}
      Clipping at $5 \times \text{median}(w_{ij})$ before KNN ranking is a
      principled consequence of Lemma~1 (the linearized formula is valid only
      for small $\|\Delta z\|$). This prevents graph oscillation (2-cycles).

    \item \textbf{Self-consistent iteration structure.}
      The trainer correctly implements: Euclidean KNN init $\to$ M-step
      (train on fixed graph) $\to$ E-step (rebuild Riemannian KNN) $\to$
      repeat.
\end{enumerate}

\subsection{Discrepancies and Concerns}

\begin{critical}[D1: Torus Data-Metric Mismatch]
The function \texttt{flat\_torus()} in \texttt{data/synthetic.py} generates:
\[
    f(\theta, \phi) = \bigl((R + r\cos\phi)\cos\theta,\;
    (R + r\cos\phi)\sin\theta,\; r\sin\phi\bigr) \in \RR^3
\]
This is the standard embedded torus, whose induced Riemannian metric is:
\[
    ds^2 = (R + r\cos\phi)^2\,d\theta^2 + r^2\,d\phi^2
\]
This metric is \textbf{not flat}: the Gaussian curvature is
$K(\phi) = \cos\phi\,/\,(r(R + r\cos\phi))$, ranging from $+1/3$ (outer ring)
to $-1$ (inner ring) for $R=2, r=1$.

However, \texttt{compute\_true\_geodesic\_distances(..., "torus")} uses:
\[
    d = \sqrt{(R\,\Delta\theta)^2 + (r\,\Delta\phi)^2}
\]
which is the metric of the \emph{abstract} flat torus
$\RR^2 / (2\pi R\ZZ \times 2\pi r\ZZ)$, NOT the induced metric of the
embedded torus.

\textbf{Impact}: The curvature proxy measures the actual curvature of the
embedded torus (nonzero), while the paper expects it to be zero. The isometry
evaluation compares learned distances against incorrect ground truth.
\end{critical}

\begin{critical}[D2: Edge KL Not in Theoretical Loss]
The implementation's loss is:
\[
    \mathcal{L} = \mathcal{L}_{\text{recon}}
    + \beta_n\,\mathcal{L}_{\text{node\_KL}}
    + \lambda\,\mathcal{L}_{\text{Riem}}
    + \beta_e\,\mathcal{L}_{\text{edge\_KL}}
\]
with $\beta_e = 10^{-3}$. The theoretical Eq.~(3) has only three terms.
The edge KL regularizer $\text{KL}[q(e_{ij} \mid \Delta z) \| \mathcal{N}(0,I)]$
acts as a bottleneck on edge codes, which may interact with the Eckart-Young
analysis of Prop.~1. Under the theoretical free-code analysis, $e_{ij}^*$
is unconstrained; the KL penalty biases $e_{ij}$ toward zero, potentially
affecting which directions $W$ captures.
\end{critical}

\begin{suggestion}[D3: Decorrelation Loss Correctly Omitted]
The decorrelation loss is defined in \texttt{loss.py} but intentionally
excluded from the total loss. The code comment correctly explains why:
the Stiefel constraint ($W^\top W = I_k$) makes the decorrelation target
$\text{Cov}(\mu_e) = I_k$ incorrect (the correct target would be
$\text{Cov}(\mu_e) = \Lambda_k$, the eigenvalues of $M_\theta$). This is a
correct architectural decision consistent with the theory.
\end{suggestion}

\begin{suggestion}[D4: Sphere Sampling Not Uniform]
The sphere data generation uses:
\begin{verbatim}
    theta = rng.uniform(0, pi, n_points)
    phi = rng.uniform(0, 2*pi, n_points)
\end{verbatim}
This is NOT uniform on $S^2$: the density is $\rho(\theta) \propto 1$ in
$\theta$, but the Riemannian volume element is $\sin\theta\,d\theta\,d\phi$.
Uniform sampling requires $\theta = \arccos(1 - 2U)$ for $U \sim
\text{Uniform}[0,1]$. The non-uniform density violates Assumption~A3 (density
bounded below). The practical impact may be small because the theory only
requires $\rho_{\min} > 0$, not exact uniformity.
\end{suggestion}


%% ====================================================================
\section{Gap Analysis}
\label{sec:gaps}
%% ====================================================================

\subsection{Gap 1: Assumption A4' -- The Metric Approximation}
\label{sec:a4}

\textbf{Status}: Open in the paper, partially resolvable.

\medskip
The paper needs:
\begin{equation}\label{eq:a4}
    \|G^*(z_i) - g(z_i)\|_{\op} \leq C_G\,\delta_{\mathrm{rec}} + O(r_n^2)
\end{equation}
but can only derive (Remark~5):
\begin{equation}\label{eq:a4_derived}
    \|G^*(z_i) - g(z_i)\|_{\op} \leq C\!\left(\frac{\delta_{\mathrm{rec}}}{r_n}
    + \kappa r_n\right)
\end{equation}
The $1/r_n$ factor is the obstacle. As $n \to \infty$, $r_n \to 0$, and
$\delta_{\mathrm{rec}}/r_n$ may diverge unless $\delta_{\mathrm{rec}}$ shrinks
faster than $r_n$.

\textbf{Resolution path}: See Section~\ref{sec:bridge_a4} for our proposed
Sobolev interpolation argument.

\subsection{Gap 2: Curvature Proxy Scaling}
\label{sec:kappa}

\textbf{Status}: Internal contradiction in the paper.

\medskip
\textbf{Claim in Eq.~(16)}: $\hat{\kappa} \to 2/R = 2\sqrt{K}$ (scale-invariant).

\textbf{Claim in Remark~8}: $\hat{\kappa} \to (2/R)\sqrt{3/G}$ (for random embedding).

\begin{proofbox}[Derivation: Why Eq.~(16) Is Incorrect for Random Embeddings]
Consider the sphere $S^2(R)$ embedded via $f = A \circ g$ where
$g: \RR^2 \to \RR^3$ is the sphere parametrization and
$A \in \RR^{G \times 3}$ has i.i.d.\ $\mathcal{N}(0, 1/\sqrt{3})$ entries.

For an equilateral triangle with edge vectors $u = \varepsilon e_1$,
$v = \varepsilon e_2$ at the north pole:

\textbf{In 3D} ($g$ alone):
$c_{ijk}^{3D} = (0, 0, 2\varepsilon^2/R)^\top$ (normal direction),
$\ell_{ij}^{3D} = (\varepsilon, 0, 0)^\top$,
$\ell_{ik}^{3D} = (0, \varepsilon, 0)^\top$.

\textbf{In $\RR^G$} ($f = Ag$): For any fixed unit vector $w \in \RR^3$,
$\|Aw\|^2 = \sum_{l=1}^G (A_l \cdot w)^2$ where $A_l \cdot w \sim
\mathcal{N}(0, 1/3)$, so $\mathbb{E}[\|Aw\|^2] = G/3$ and by concentration
$\|Aw\| \approx \sqrt{G/3}$ for all unit $w$.

Therefore:
\begin{align*}
    \text{Numerator:}\quad \|c_{ijk}\|_{\RR^G}
      &= \|A c_{ijk}^{3D}\| \approx \sqrt{G/3}\;\|c_{ijk}^{3D}\|
      = \sqrt{G/3}\cdot\frac{2\varepsilon^2}{R} \\
    \text{Denominator:}\quad \|\ell_{ij}\|_{\RR^G}\,\|\ell_{ik}\|_{\RR^G}
      &\approx \left(\sqrt{G/3}\,\varepsilon\right)^2
      = \frac{G}{3}\,\varepsilon^2
\end{align*}
\begin{equation}
    \hat{\kappa}
    = \frac{\sqrt{G/3}\cdot 2\varepsilon^2/R}{(G/3)\,\varepsilon^2}
    = \frac{2}{R}\cdot\frac{1}{\sqrt{G/3}}
    = \frac{2}{R}\sqrt{\frac{3}{G}}
\end{equation}
For $R = 1$, $G = 50$: $\hat{\kappa}^* = 2\sqrt{3/50} = 0.4899$.

\medskip
\textbf{The error in Eq.~(16)}: The claim ``scale-invariant'' is correct for
rescaling $z$ (latent coordinates) but \textbf{incorrect for the ambient
embedding}. The denominator is a product of two ambient-space norms, each
scaling as $\sqrt{G/3}$, while the numerator scales as $\sqrt{G/3}$ (not
$(G/3)$). The net effect is a $1/\sqrt{G/3}$ factor that Eq.~(16) omits.

\medskip
\textbf{Consequence for experiments}: The paper compares measured
$\hat{\kappa} = 1.442$ against the prediction $2.000$. The correct prediction
for the random-embedding setup is $0.4899$. Under this correction, the
measured value $1.442$ overshoots by a factor of $\approx 3$, which may
indicate that the decoder has not fully converged or that finite-triangle-size
effects dominate.
\end{proofbox}

\subsection{Gap 3: The Torus Is Not Flat}
\label{sec:torus}

\textbf{Status}: Critical experimental error.

\begin{proofbox}[The Embedded Torus Has Nonzero Curvature]
The standard parametric torus in $\RR^3$ is:
\[
    f(\theta, \phi) = \bigl((R+r\cos\phi)\cos\theta,\;
    (R+r\cos\phi)\sin\theta,\; r\sin\phi\bigr)
\]
Its induced metric tensor is:
\[
    g =
    \begin{pmatrix}
        (R+r\cos\phi)^2 & 0 \\
        0 & r^2
    \end{pmatrix}
\]
The Gaussian curvature is:
\begin{equation}
    K(\theta,\phi) = \frac{\cos\phi}{r(R + r\cos\phi)}
\end{equation}
For $R = 2$, $r = 1$:
\[
    K_{\max} = K(0) = \frac{1}{1\cdot(2+1)} = \frac{1}{3},
    \qquad
    K_{\min} = K(\pi) = \frac{-1}{1\cdot(2-1)} = -1
\]
The curvature ranges from $-1$ to $+1/3$ across the torus surface.

By the Gauss-Bonnet theorem, $\int_M K\,dA = 2\pi\chi(T^2) = 0$, so the
total curvature integrates to zero, but the pointwise curvature is
\textbf{nowhere zero} (except at $\phi = \pi/2$ and $\phi = 3\pi/2$).
\end{proofbox}

\textbf{The flat torus} is the abstract quotient $\RR^2/(2\pi R\ZZ \times
2\pi r\ZZ)$, which cannot be isometrically embedded in $\RR^3$ by a theorem
of Weyl (smooth isometric immersions of flat tori in $\RR^3$ do not exist
globally; the Nash-Kuiper $C^1$ embedding is highly fractal).

The \textbf{Clifford torus} provides a smooth isometric embedding in $\RR^4$:
\begin{equation}\label{eq:clifford}
    f_{\text{Cliff}}(\theta, \phi)
    = (R\cos\theta,\; R\sin\theta,\; r\cos\phi,\; r\sin\phi) \in \RR^4
\end{equation}
This has:
\begin{itemize}
    \item Induced metric: $ds^2 = R^2\,d\theta^2 + r^2\,d\phi^2$ (constant
      coefficients = flat).
    \item Gaussian curvature: $K = 0$ everywhere.
    \item The geodesic distance formula
      $d = \sqrt{(R\,\Delta\theta)^2 + (r\,\Delta\phi)^2}$ (with
      $\Delta\theta, \Delta\phi$ mod $2\pi$) is \textbf{exact} for this embedding.
\end{itemize}

\subsection{Gap 4: Global Convergence and Basin of Attraction}

\textbf{Status}: Open, acknowledged by the paper.

The convergence proof (Prop.~3) is local: it requires initialization near the
fixed point. The Euclidean KNN graph used for initialization may or may not
lie in the basin of attraction. For manifolds with high curvature or complex
topology, the Euclidean graph can be far from the Riemannian graph, and the
initial decoder may converge to a suboptimal local minimum.

No formal characterization of the basin of attraction is provided. This is
inherent to non-convex optimization and difficult to resolve in full
generality, but conditions of the form ``$K_{\max} r_n^{(0)} < c$ implies
the Euclidean KNN equals the Riemannian KNN'' would strengthen the result.


%% ====================================================================
\section{Bridging the Gaps: New Mathematical Results}
\label{sec:bridge}
%% ====================================================================

\subsection{Bridging A4': Sobolev Interpolation for Neural Decoders}
\label{sec:bridge_a4}

We provide a partial resolution of the Assumption~A4' gap.

\begin{lemma}[Jacobian control via Sobolev interpolation]
\label{lem:sobolev}
Let $f_\theta, f_{\mathrm{true}}: \RR^d \to \RR^G$ be $C^2$ functions with
uniformly bounded Hessians:
\[
    \sup_z \|H_{f_\theta}(z)\|_{\op} \leq \kappa_\theta, \qquad
    \sup_z \|H_{f_{\mathrm{true}}}(z)\|_{\op} \leq \kappa_{\mathrm{true}}
\]
Let $\{z_i\}_{i=1}^n$ be an $r_n$-covering of a compact domain
$\Omega \subset \RR^d$ (every point in $\Omega$ is within $r_n$ of some
$z_i$). Define $\delta = \max_i \|f_\theta(z_i) - f_{\mathrm{true}}(z_i)\|$
and $\kappa = \kappa_\theta + \kappa_{\mathrm{true}}$.

Then for any $z$ in the convex hull of $\{z_i\}$:
\begin{equation}\label{eq:jac_bound}
    \|J_{f_\theta}(z) - J_{f_{\mathrm{true}}}(z)\|_{\op}
    \leq \frac{2\delta}{r_n} + \kappa\,r_n
\end{equation}
\end{lemma}

\begin{proof}
For any training point $z_i$ and direction $v$ with $\|v\| = 1$, set
$z_j = z_i + r_n v$ (which is within $\Omega$ or its convex hull for
sufficiently dense sampling). By Taylor expansion:
\begin{align*}
    f_\theta(z_j) - f_\theta(z_i)
    &= J_{f_\theta}(z_i)\,(z_j - z_i) + O(\kappa_\theta\,r_n^2) \\
    f_{\mathrm{true}}(z_j) - f_{\mathrm{true}}(z_i)
    &= J_{f_{\mathrm{true}}}(z_i)\,(z_j - z_i) + O(\kappa_{\mathrm{true}}\,r_n^2)
\end{align*}
Subtracting:
\[
    (J_{f_\theta}(z_i) - J_{f_{\mathrm{true}}}(z_i))\,r_n v
    = \bigl(f_\theta(z_j) - f_{\mathrm{true}}(z_j)\bigr)
    - \bigl(f_\theta(z_i) - f_{\mathrm{true}}(z_i)\bigr)
    + O(\kappa\,r_n^2)
\]
Taking norms and using $\|f_\theta(z_l) - f_{\mathrm{true}}(z_l)\| \leq \delta$
for training points:
\[
    r_n\,\|(J_{f_\theta}(z_i) - J_{f_{\mathrm{true}}}(z_i))v\|
    \leq 2\delta + \kappa\,r_n^2
\]
Dividing by $r_n$ and maximizing over $\|v\| = 1$:
$\|J_{f_\theta}(z_i) - J_{f_{\mathrm{true}}}(z_i)\|_{\op} \leq 2\delta/r_n
+ \kappa\,r_n$.
For non-training points $z$ with $\|z - z_i\| \leq r_n$, the Hessian bound
gives $\|J_f(z) - J_f(z_i)\|_{\op} \leq \kappa r_n$, which is absorbed.
\end{proof}

\begin{corollary}[Sufficient condition for A4']
\label{cor:a4}
If the decoder reconstruction error satisfies
$\delta_{\mathrm{rec}} = O(r_n^2)$ -- i.e., the decoder has sufficient
capacity to achieve reconstruction accuracy scaling as the squared KNN
radius -- then:
\begin{equation}
    \|G^*(z_i) - g(z_i)\|_{\op}
    \leq C\!\left(\frac{r_n^2}{r_n} + \kappa r_n\right)
    = O(r_n)
\end{equation}
and Assumption~A4' holds automatically with the
$C_G\delta_{\mathrm{rec}} + O(r_n^2)$ form replaced by $O(r_n)$.
In this regime, the $C_2\,\delta_{\mathrm{rec}}$ term in Theorem~2
is absorbed into $C_1\,r_n$, and the isometry bound simplifies to:
\begin{equation}
    \mathbb{E}_{i,j}\!\left[
    |\dR^*(z_i,z_j) - \dM(x_i,x_j)|\right]
    \leq C_1'\,r_n + \frac{K_{\max}}{6}\,r_n^3 + C_3\,e^{-\lambda T/2}
\end{equation}
\end{corollary}

\begin{remark}[When does $\delta_{\mathrm{rec}} = O(r_n^2)$ hold?]
For a neural decoder with $L$ hidden layers of width $m$ and SiLU activation:
\begin{enumerate}[label=(\alph*)]
    \item \textbf{Universal approximation}: By the Barron-type theorem for
      smooth functions, for any $C^2$ target function on a compact domain,
      a two-layer network of width $m$ achieves
      $\delta_{\mathrm{rec}} = O(1/m)$ (Barron 1993).
    \item \textbf{Scaling requirement}: $\delta_{\mathrm{rec}} = O(r_n^2)
      = O((\log n / n)^{2/d})$ is satisfied when
      $m = \Omega((n/\log n)^{2/d})$.
    \item For $d = 2$: $m = \Omega(n / \log n)$. For $n = 3000$:
      $m \gtrsim 400$. The decoder architecture (hidden widths 128, 256) may
      be marginal.
\end{enumerate}
\end{remark}

\subsection{Bridging the Curvature Proxy: Corrected Scaling Law}

\begin{proposition}[Corrected curvature proxy for random embeddings]
\label{prop:kappa_corrected}
Let $M^d \hookrightarrow \RR^{d+1}$ be a smooth submanifold with sectional
curvature $K$ at a point $p$. Let $A \in \RR^{G \times (d+1)}$ be a random
embedding with i.i.d.\ $\mathcal{N}(0, \sigma^2)$ entries. Define the
curvature proxy $\hat{\kappa}$ using the ambient-space norms as in the paper.

Then as the triangle scale $\varepsilon \to 0$:
\begin{equation}\label{eq:kappa_correct}
    \hat{\kappa}
    \;\to\;
    \frac{\hat{\kappa}^{d+1}}{\sigma\sqrt{G}}
    \;=\;
    \frac{2\sqrt{|K|}}{\sigma\sqrt{G}}
\end{equation}
where $\hat{\kappa}^{d+1}$ is the curvature proxy computed in the
$(d+1)$-dimensional pre-embedding space.

For the experimental setup ($R=1$, $G=50$, $\sigma = 1/\sqrt{3}$):
\begin{equation}
    \hat{\kappa}^*_{\text{sphere}}
    = \frac{2}{\sqrt{50/3}} = 2\sqrt{\frac{3}{50}} = 0.4899
\end{equation}
For the flat torus with Clifford embedding:
$\hat{\kappa}^*_{\text{flat torus}} = 0$.
\end{proposition}

\begin{proof}
For any unit vector $w \in \RR^{d+1}$, the random projection satisfies
$\|Aw\|^2 = \sigma^2 \sum_{g=1}^G Z_g^2$ where $Z_g \sim \mathcal{N}(0,1)$.
By the law of large numbers, $\|Aw\|/(\sigma\sqrt{G}) \to 1$ as $G \to \infty$.

The closure vector $c \in \RR^{d+1}$ and log maps $\ell \in \RR^{d+1}$
satisfy $\|Ac\| \approx \sigma\sqrt{G}\,\|c\|$ and
$\|A\ell\| \approx \sigma\sqrt{G}\,\|\ell\|$.

The curvature proxy ratio:
\[
    \hat{\kappa} = \frac{\|Ac\|}{\|A\ell_{ij}\|\,\|A\ell_{ik}\|}
    \approx \frac{\sigma\sqrt{G}\,\|c\|}{(\sigma\sqrt{G})^2\,\|\ell_{ij}\|\,\|\ell_{ik}\|}
    = \frac{\|c\|}{\sigma\sqrt{G}\,\|\ell_{ij}\|\,\|\ell_{ik}\|}
    = \frac{\hat{\kappa}^{d+1}}{\sigma\sqrt{G}}
\]
Substituting $\hat{\kappa}^{d+1} = 2/R = 2\sqrt{K}$ for the sphere and
$\sigma = 1/\sqrt{3}$:
$\hat{\kappa} = 2\sqrt{K} / (\sqrt{G/3}) = 2\sqrt{3K/G}$.
\end{proof}


%% ====================================================================
\section{Analysis: Sphere vs.\ Torus Performance Gap}
\label{sec:sphere_vs_torus}
%% ====================================================================

The paper reports that the sphere benefits substantially more from Riemannian
training (34\% curvature improvement) than the torus (13\%), and identifies
three explanations: topology, linearization outliers, and the isometry floor.
Our analysis reveals that the dominant factor is simpler and more fundamental.

\subsection{The Primary Cause: Wrong Manifold}

\begin{critical}[Root Cause of Torus Underperformance]
The torus experiment tests on the \textbf{standard embedded torus} in $\RR^3$,
which has Gaussian curvature $K(\phi) = \cos\phi/(r(R+r\cos\phi))$, but
evaluates against the \textbf{flat torus metric} $d = \sqrt{(R\,d\theta)^2 +
(r\,d\phi)^2}$. This creates three compounding errors:

\begin{enumerate}
    \item \textbf{Curvature proxy}: The measured $\hat{\kappa} = 2.212$
      reflects the \emph{actual} nonzero curvature of the embedded torus, not
      a failure of the method. The theoretical prediction should be based on
      the curvature of the embedded torus, not zero.

    \item \textbf{Isometry evaluation}: Comparing learned Riemannian distances
      against the flat-torus geodesics introduces systematic error. On the
      outer ring ($\phi \approx 0$), the true embedded metric gives
      $(R+r)^2 d\theta^2 + r^2 d\phi^2$ while the flat metric uses
      $R^2 d\theta^2 + r^2 d\phi^2$. The relative difference in the
      $\theta$-direction is $(R+r)^2/R^2 - 1 = (2r/R + r^2/R^2)$, which
      equals $125\%$ for $R=2, r=1$.

    \item \textbf{Graph construction}: The Riemannian KNN graph correctly
      uses the embedded metric (via JVP), but the evaluation assumes a
      different metric. The graph may be optimal for the true geometry while
      appearing suboptimal under the flat evaluation.
\end{enumerate}
\end{critical}

\subsection{Secondary Causes (Genuine, After Fixing the Manifold)}

Even with the correct flat torus (Clifford embedding), several genuine
challenges remain:

\paragraph{(a) Topological periodicity.}
The flat torus has fundamental group $\ZZ \times \ZZ$. An encoder mapping
$T^2 \to \RR^2$ must ``unwrap'' both periodic directions. This unwrapping
requires the encoder to learn a discontinuous map (cutting the torus along
two circles), which a smooth neural network approximates poorly at the
cuts. This causes ``boundary effects'' where points near the cut have
distorted latent neighborhoods.

The sphere $S^2$ is simply connected, so no unwrapping is needed (except at
the antipodal point, which is a single point of measure zero).

\paragraph{(b) Latent space mismatch.}
A Euclidean latent space $\RR^2$ cannot globally represent a torus without
distortion -- this is a topological obstruction (the torus is compact but
$\RR^2$ is not). The Euclidean KL prior $\mathcal{N}(0, I_2)$ further biases
the latent representation toward a single Gaussian mode, conflicting with the
torus's periodic structure.

For the sphere, the same obstruction exists (the sphere is also compact), but
the single connected component is easier to represent in a single Gaussian
mode.

\paragraph{(c) Optimizer state management.}
As the paper notes, resetting Adam's momentum helps on the sphere but hurts
on the torus. This is because the torus requires learning two independent
periodic directions simultaneously, which benefits from accumulated momentum
across iterations.

\subsection{Mathematical Characterization of the Torus Isometry Floor}

\begin{proposition}[Isometry lower bound for torus in Euclidean latent space]
\label{prop:torus_floor}
Let $f_\theta: \RR^2 \to \RR^G$ be any $C^1$ decoder and
$g_\phi: \RR^G \to \RR^2$ any encoder. For the flat torus $T^2$ with
circumferences $2\pi R$ and $2\pi r$, the worst-case isometry error satisfies:
\begin{equation}
    \max_{x_i, x_j \in T^2}
    |\dR^*(z_i, z_j) - \dM(x_i, x_j)|
    \geq \min(\pi R, \pi r)
\end{equation}
\end{proposition}

\begin{proof}[Proof sketch]
Consider the shortest geodesic loop $\gamma$ around the smaller circle of the
torus (length $2\pi r$). The encoder maps this loop to a closed curve in
$\RR^2$. By the intermediate value theorem, there exist two points on this
curve that map to the same latent point (the curve must cross itself in
$\RR^2$, since a non-self-intersecting closed curve in $\RR^2$ encloses a
disk, but the torus is not a disk).

At the crossing: $z_i = z_j$ but $\dM(x_i, x_j) > 0$. The Riemannian
distance $\dR^*(z_i, z_j) = 0$ while $\dM(x_i, x_j) \leq \pi r$ (the
crossing is at most halfway around the loop). Hence the error is at least
$\pi r$.

More precisely: for any continuous $g: T^2 \to \RR^2$, the Borsuk-Ulam-type
argument shows that there exist antipodal points on each generating circle
that are mapped to the same point in $\RR^2$, giving error $\geq \pi r$.
\end{proof}

\begin{remark}
This lower bound is $\pi r = \pi \approx 3.14$ for $r=1$. The measured
Euclidean latent MAE is $\approx 1.35$, which is below this bound because
the MAE averages over all pairs (not worst-case) and uses random pair
sampling that avoids the worst configurations.
\end{remark}


%% ====================================================================
\section{Recommendations and Action Items}
\label{sec:recommendations}
%% ====================================================================

\subsection{Critical Fixes (Must-Do Before Submission)}

\begin{enumerate}[label=\textbf{F\arabic*.}]
    \item \textbf{Fix the torus experiment.}
      Replace the standard embedded torus in $\RR^3$ with the Clifford torus
      in $\RR^4$ (Eq.~\ref{eq:clifford}). This requires:
      \begin{itemize}
          \item Modify \texttt{flat\_torus()} to generate
            $f(\theta,\phi) = (R\cos\theta, R\sin\theta, r\cos\phi, r\sin\phi)$.
          \item Change the embedding matrix to $A \in \RR^{G \times 4}$.
          \item The geodesic distance formula in
            \texttt{compute\_true\_geodesic\_distances} is already correct for
            the flat metric.
      \end{itemize}
      \textbf{Expected impact}: The curvature proxy should approach zero,
      validating the negative test. The isometry evaluation will use the
      correct ground truth.

    \item \textbf{Resolve the curvature proxy inconsistency.}
      Remove or correct Eq.~(16). The correct prediction for random embeddings
      is Remark~8's formula: $\hat{\kappa} = (2/R)\sqrt{3/G}$. Update the
      experiments section to compare against this value ($0.4899$ for $R=1,
      G=50$).

    \item \textbf{Add the edge KL term to the theoretical loss.}
      Either justify its absence from the theory (e.g., show that for
      $\beta_e \ll 1$ the effect is perturbative) or incorporate it into the
      convergence analysis.
\end{enumerate}

\subsection{Important Improvements}

\begin{enumerate}[label=\textbf{I\arabic*.}]
    \item \textbf{Strengthen the presentation of A4'.}
      Add Corollary~\ref{cor:a4} showing that $\delta_{\mathrm{rec}} =
      O(r_n^2)$ implies A4'. Discuss the decoder capacity requirements:
      for $d=2$, $n=3000$, hidden width $\geq 400$ may be needed.

    \item \textbf{Fix the sphere sampling.}
      Use $\theta = \arccos(1 - 2U)$ for truly uniform sampling on $S^2$.
      This ensures $\rho_{\min} > 0$ uniformly, satisfying Assumption~A3
      without relying on the non-uniform density being bounded below.

    \item \textbf{Provide an additional manifold experiment.}
      The Swiss roll (already implemented in \texttt{data/synthetic.py}) would
      be a good third test case: it has varying curvature, no topological
      obstruction (it is homeomorphic to $\RR^2$), and is a standard benchmark
      for manifold learning methods.

    \item \textbf{Empirical verification of A4'.}
      At convergence, compute $\|G^*(z_i) - g(z_i)\|_{\op}$ for the sphere
      (where the true metric is known analytically) and plot against $r_n$.
      If the scaling is $O(r_n)$ rather than $O(\delta_{\mathrm{rec}}/r_n)$,
      this provides empirical evidence that A4' holds in the overparameterized
      regime.

    \item \textbf{Multi-scale self-consistent iteration for the torus.}
      Start with a large neighborhood radius (capturing global topology) and
      gradually reduce it (capturing local geometry). This ``coarse-to-fine''
      strategy may help the torus converge faster by first resolving the
      periodic structure and then refining local distances.
\end{enumerate}

\subsection{Nice-to-Have Enhancements}

\begin{enumerate}[label=\textbf{N\arabic*.}]
    \item \textbf{Spectral gap and convergence rate.}
      Open Problem~1 asks about the relationship between $\lambda$ and the
      manifold's spectral gap $\delta$. For the Laplace-Beltrami operator on
      $S^2$, $\delta = 2/R^2$. On $T^2$, $\delta = \min(1/R^2, 1/r^2)$. A
      formal connection would explain the observed convergence speed difference
      ($\sim15$ vs.\ $\sim98$ iterations).

    \item \textbf{Curvature-aware training.}
      Use the curvature proxy $\hat{\kappa}$ as a training signal: regions with
      high $\hat{\kappa}$ need denser neighborhoods or stronger Riemannian
      regularization. This could be implemented as an adaptive $\lambda$
      weight or non-uniform KNN ($k_i$ varying by local curvature).

    \item \textbf{Topologically-aware latent spaces.}
      For the torus, replace $\RR^2$ with a periodic latent space (e.g.,
      $[0, 2\pi)^2$ with periodic boundary conditions). This eliminates the
      topological obstruction of Proposition~\ref{prop:torus_floor} and should
      dramatically improve torus performance.
\end{enumerate}


%% ====================================================================
\section{Summary of Findings}
\label{sec:summary}
%% ====================================================================

\begin{center}
\small
\begin{tabular}{p{4.5cm}p{2cm}p{7cm}}
\hline
\textbf{Finding} & \textbf{Severity} & \textbf{Resolution} \\
\hline
Torus uses non-flat embedding &
  Critical &
  Use Clifford torus in $\RR^4$ \\
Curvature proxy Eq.~(16) vs.\ Remark~8 &
  Critical &
  Correct Eq.~(16) to $(2/R)\sqrt{3/G}$ \\
Assumption A4' unresolved &
  Major &
  Corollary~\ref{cor:a4}: suffices when $\delta_{\mathrm{rec}} = O(r_n^2)$ \\
Edge KL not in theoretical loss &
  Moderate &
  Add to Eq.~(3) or justify perturbative regime \\
Sphere sampling non-uniform &
  Minor &
  Use $\theta = \arccos(1-2U)$ \\
No global convergence &
  Major (open) &
  Characterize basin of attraction \\
Torus convergence slow (98 iters) &
  Moderate &
  Multi-scale iteration + periodic latent space \\
\hline
\end{tabular}
\end{center}

\subsection{Final Verdict}

The SCR-VAE introduces a genuinely novel framework for learning isometric
generative models. The self-consistent iteration, the fixed-point
characterization, and the curvature certificate are original contributions
that advance the field. However, the experimental validation has a critical
flaw (wrong torus manifold) and the curvature proxy's theoretical prediction
contains an internal contradiction that affects the interpretation of all
curvature experiments.

The core theory is sound: Lemma~1, Theorem~1, Proposition~3, and the
convergence analysis (Proposition~3/Corollary~1) are all correct. The main
theoretical gap (Assumption~A4') is partially bridgeable via the Sobolev
interpolation argument, and fully resolvable in the high-capacity decoder
regime ($\delta_{\mathrm{rec}} = O(r_n^2)$).

With the three critical fixes (torus manifold, curvature scaling, loss
alignment), the paper would present a strong contribution to NeurIPS.

\end{document}
